% =====================================================================
%  figures/binomial-barras.tex
%  Distribución binomial (n=10, p=0,324).
%  (Reemplaza a img/barras.png)
% =====================================================================
\begin{figure}[H]
\centering
\begin{tikzpicture}
\begin{axis}[
  ybar, bar width=20pt,
  title={Distribución binomial ($n=10$, $p=0{,}324$)},
  xlabel={Número de estudiantes con asistencia perfecta ($k$)},
  ylabel={Probabilidad (\%)},
  ymin=0, ymax=30,
  xtick=data,
  enlarge x limits=0.08,
  clip=false,
]
% Todas las barras en azul
\addplot+[fill=azul, draw=azul] coordinates {
  (0,1.98) (1,9.49) (2,20.67) (3,26.33) (4,22.07) (5,12.70)
  (6,5.06) (7,1.39) (8,0.25) (9,0.03) (10,0.00)
};
% Barra k=3 resaltada en naranja (rectángulo superpuesto, cubre la azul)
\fill[naranja, draw=naranja] (axis cs:2.6,0) rectangle (axis cs:3.4,28);
\end{axis}
\end{tikzpicture}
\caption{Gráfico de distribución binomial.}
\label{fig:grafico_probabilidad}
\end{figure}
